% =====================================================================
%  terms.tex —— 英汉术语对照（由 tools/make_terms_appendix.py 从 术语对照表.tsv 生成，勿手改）
%  共 9 组 / 593 条
% =====================================================================
\chapter*{附录：英汉术语对照}
\addcontentsline{toc}{chapter}{附录：英汉术语对照}
\markboth{附录：英汉术语对照}{附录：英汉术语对照}
\begingroup
\small
下表是本书译名与英文原词的一一对应。凡表中已有的词，译文中一律照用；表中未列的，按各章首次出现处的括注为准。

\filbreak
\vspace{0.9em}
\noindent\textbf{一、量子力学基本概念}\par\nobreak\vspace{0.35em}
\begin{multicols}{2}
\raggedright
\setlength{\parindent}{0pt}
\setlength{\parskip}{0.16em}
\linespread{1.02}\selectfont
quantum mechanics\hspace{0.9em}量子力学\par
state\hspace{0.9em}态\par
state vector\hspace{0.9em}态向量\par
state space\hspace{0.9em}状态空间\par
Hilbert space\hspace{0.9em}希尔伯特空间\par
vector\hspace{0.9em}向量\par
basis\hspace{0.9em}基\par
orthonormal basis\hspace{0.9em}正交归一基\par
orthogonal basis\hspace{0.9em}正交基\par
completeness\hspace{0.9em}完备性\par
complete\hspace{0.9em}完备的\par
orthogonal\hspace{0.9em}正交的\par
normalized\hspace{0.9em}归一的\par
normalization\hspace{0.9em}归一化\par
inner product\hspace{0.9em}内积\par
norm\hspace{0.9em}范数\par
superposition\hspace{0.9em}叠加\par
linear combination\hspace{0.9em}线性组合\par
wave function\hspace{0.9em}波函数\par
wave packet\hspace{0.9em}波包\par
probability amplitude\hspace{0.9em}概率振幅\par
probability density\hspace{0.9em}概率密度\par
probability distribution\hspace{0.9em}概率分布\par
expectation value\hspace{0.9em}期望值\par
standard deviation\hspace{0.9em}标准差\par
uncertainty relation\hspace{0.9em}不确定性关系\par
uncertainty principle\hspace{0.9em}不确定性原理\par
measurement\hspace{0.9em}测量\par
observable\hspace{0.9em}可观测量\par
observation\hspace{0.9em}观测\par
collapse\hspace{0.9em}塌缩\par
probabilistic interpretation\hspace{0.9em}概率诠释\par
Born rule\hspace{0.9em}玻恩规则\par
eigenstate\hspace{0.9em}本征态\par
eigenvalue\hspace{0.9em}本征值\par
eigenvector\hspace{0.9em}本征向量\par
eigenfunction\hspace{0.9em}本征函数\par
eigenspace\hspace{0.9em}本征子空间\par
degeneracy\hspace{0.9em}简并\par
degenerate\hspace{0.9em}简并的\par
spectrum\hspace{0.9em}谱\par
discrete spectrum\hspace{0.9em}离散谱\par
continuous spectrum\hspace{0.9em}连续谱\par
point spectrum\hspace{0.9em}点谱\par
entanglement\hspace{0.9em}纠缠\par
entangled state\hspace{0.9em}纠缠态\par
product state\hspace{0.9em}积态\par
density matrix\hspace{0.9em}密度矩阵\par
tensor product\hspace{0.9em}张量积\par
tensor product space\hspace{0.9em}张量积空间\par
quantum harmonic oscillator\hspace{0.9em}量子谐振子\par
harmonic oscillator\hspace{0.9em}谐振子\par
free particle\hspace{0.9em}自由粒子\par
square well\hspace{0.9em}方势阱\par
potential\hspace{0.9em}势\par
potential well\hspace{0.9em}势阱\par
barrier\hspace{0.9em}势垒\par
bound state\hspace{0.9em}束缚态\par
scattering state\hspace{0.9em}散射态\par
scattering\hspace{0.9em}散射\par
cross section\hspace{0.9em}截面\par
transmission coefficient\hspace{0.9em}透射系数\par
reflection coefficient\hspace{0.9em}反射系数\par
boundary condition\hspace{0.9em}边界条件\par
matching condition\hspace{0.9em}衔接条件\par
asymptotic\hspace{0.9em}渐近的\par
asymptotic behavior\hspace{0.9em}渐近行为\par
turning point\hspace{0.9em}转折点\par
WKB approximation\hspace{0.9em}WKB 近似\par
semi-classical\hspace{0.9em}半经典的\par
semi-classical approximation\hspace{0.9em}半经典近似\par
Bohr–Sommerfeld\hspace{0.9em}玻尔–索末菲\par
Kepler problem\hspace{0.9em}开普勒问题\par
Coulomb potential\hspace{0.9em}库仑势\par
Runge–Lenz vector\hspace{0.9em}龙格–楞次向量\par
hydrogen atom\hspace{0.9em}氢原子\par
symmetry\hspace{0.9em}对称性\par
conservation law\hspace{0.9em}守恒律\par
angular momentum\hspace{0.9em}角动量\par
spin\hspace{0.9em}自旋\par
parity\hspace{0.9em}宇称\par
time reversal\hspace{0.9em}时间反演\par
translation\hspace{0.9em}平移\par
generator\hspace{0.9em}生成元\par
representation\hspace{0.9em}表示\par
unitary representation\hspace{0.9em}酉表示\par
irreducible representation\hspace{0.9em}不可约表示\par
Lie algebra\hspace{0.9em}李代数\par
Lie group\hspace{0.9em}李群\par
commutator\hspace{0.9em}交换子\par
canonical commutation relation\hspace{0.9em}正则对易关系\par
commutation relation\hspace{0.9em}对易关系\par
Hamiltonian\hspace{0.9em}哈密顿量\par
Lagrangian\hspace{0.9em}拉格朗日量\par
action\hspace{0.9em}作用量\par
conjugate momentum\hspace{0.9em}共轭动量\par
Poisson bracket\hspace{0.9em}泊松括号\par
phase space\hspace{0.9em}相空间\par
trajectory\hspace{0.9em}轨迹\par
classical limit\hspace{0.9em}经典极限\par
quantum correction\hspace{0.9em}量子修正\par
Planck constant\hspace{0.9em}普朗克常数\par
\end{multicols}

\filbreak
\vspace{0.9em}
\noindent\textbf{二、希尔伯特空间与算子（泛函分析）}\par\nobreak\vspace{0.35em}
\begin{multicols}{2}
\raggedright
\setlength{\parindent}{0pt}
\setlength{\parskip}{0.16em}
\linespread{1.02}\selectfont
linear operator\hspace{0.9em}线性算子\par
operator\hspace{0.9em}算子\par
bounded operator\hspace{0.9em}有界算子\par
unbounded operator\hspace{0.9em}无界算子\par
self-adjoint\hspace{0.9em}自伴的\par
self-adjoint operator\hspace{0.9em}自伴算子\par
symmetric operator\hspace{0.9em}对称算子\par
Hermitian\hspace{0.9em}厄米的\par
essentially self-adjoint\hspace{0.9em}本质自伴的\par
adjoint\hspace{0.9em}伴随\par
unitary\hspace{0.9em}酉的\par
unitary operator\hspace{0.9em}酉算子\par
unitary group\hspace{0.9em}酉群\par
projection\hspace{0.9em}投影\par
orthogonal projection\hspace{0.9em}正交投影\par
resolution of the identity\hspace{0.9em}恒等分解\par
projection-valued measure\hspace{0.9em}投影值测度\par
spectral measure\hspace{0.9em}谱测度\par
Borel\hspace{0.9em}博雷尔\par
Borel functional calculus\hspace{0.9em}博雷尔泛函演算\par
spectral theorem\hspace{0.9em}谱定理\par
functional calculus\hspace{0.9em}泛函演算\par
resolvent\hspace{0.9em}预解式\par
resolvent set\hspace{0.9em}预解集\par
domain\hspace{0.9em}定义域\par
range\hspace{0.9em}值域\par
kernel\hspace{0.9em}核\par
null space\hspace{0.9em}零空间\par
graph\hspace{0.9em}图像\par
closed operator\hspace{0.9em}闭算子\par
closed graph theorem\hspace{0.9em}闭图像定理\par
densely defined\hspace{0.9em}稠密定义的\par
extension\hspace{0.9em}扩张\par
closure\hspace{0.9em}闭包\par
compact operator\hspace{0.9em}紧算子\par
rank\hspace{0.9em}秩\par
trace\hspace{0.9em}迹\par
Fredholm operator\hspace{0.9em}弗雷德霍姆算子\par
index\hspace{0.9em}指标\par
cokernel\hspace{0.9em}余核\par
finite-dimensional\hspace{0.9em}有限维的\par
infinite-dimensional\hspace{0.9em}无穷维的\par
strongly continuous\hspace{0.9em}强连续的\par
one-parameter unitary group\hspace{0.9em}单参数酉群\par
Stone's theorem\hspace{0.9em}斯通定理\par
Stone–von Neumann theorem\hspace{0.9em}斯通–冯·诺依曼定理\par
irreducibility\hspace{0.9em}不可约性\par
Schrödinger representation\hspace{0.9em}薛定谔表示\par
Fock representation\hspace{0.9em}福克表示\par
Schrödinger operator\hspace{0.9em}薛定谔算子\par
Laplacian\hspace{0.9em}拉普拉斯算子\par
essential spectrum\hspace{0.9em}本质谱\par
bra\hspace{0.9em}左矢\par
ket\hspace{0.9em}右矢\par
Dirac delta function\hspace{0.9em}狄拉克 $\delta$ 函数\par
distribution\hspace{0.9em}分布\par
tempered distribution\hspace{0.9em}缓增分布\par
test function\hspace{0.9em}检验函数\par
Fourier transform\hspace{0.9em}傅里叶变换\par
inverse Fourier transform\hspace{0.9em}傅里叶逆变换\par
convolution\hspace{0.9em}卷积\par
Parseval\hspace{0.9em}帕塞瓦尔\par
Plancherel\hspace{0.9em}普朗歇尔\par
\end{multicols}

\filbreak
\vspace{0.9em}
\noindent\textbf{三、路径积分}\par\nobreak\vspace{0.35em}
\begin{multicols}{2}
\raggedright
\setlength{\parindent}{0pt}
\setlength{\parskip}{0.16em}
\linespread{1.02}\selectfont
path integral\hspace{0.9em}路径积分\par
path integral formalism\hspace{0.9em}路径积分形式\par
propagator\hspace{0.9em}传播子\par
transition amplitude\hspace{0.9em}跃迁振幅\par
time slicing\hspace{0.9em}时间切片\par
classical path\hspace{0.9em}经典路径\par
fluctuation\hspace{0.9em}涨落\par
Gaussian integral\hspace{0.9em}高斯积分\par
Gaussian\hspace{0.9em}高斯的\par
measure\hspace{0.9em}测度\par
Wiener measure\hspace{0.9em}维纳测度\par
Euclidean\hspace{0.9em}欧氏的\par
analytic continuation\hspace{0.9em}解析延拓\par
Wick rotation\hspace{0.9em}威克转动\par
stationary phase\hspace{0.9em}驻定相位\par
stationary phase approximation\hspace{0.9em}驻定相位近似\par
stationary phase method\hspace{0.9em}驻定相位方法\par
steepest descent\hspace{0.9em}最速下降\par
saddle point\hspace{0.9em}鞍点\par
oscillatory integral\hspace{0.9em}振荡积分\par
phase function\hspace{0.9em}相位函数\par
asymptotic expansion\hspace{0.9em}渐近展开\par
Maslov index\hspace{0.9em}马斯洛夫指标\par
caustic\hspace{0.9em}焦散\par
Stokes phenomenon\hspace{0.9em}斯托克斯现象\par
Stokes line\hspace{0.9em}斯托克斯线\par
connection formula\hspace{0.9em}连接公式\par
Airy function\hspace{0.9em}艾里函数\par
Green's function\hspace{0.9em}格林函数\par
semi-classical Green's function\hspace{0.9em}半经典格林函数\par
perturbation\hspace{0.9em}微扰\par
perturbative expansion\hspace{0.9em}微扰展开\par
Feynman diagram\hspace{0.9em}费曼图\par
connected diagram\hspace{0.9em}连通图\par
vacuum\hspace{0.9em}真空\par
partition function\hspace{0.9em}配分函数\par
scattering matrix\hspace{0.9em}散射矩阵\par
annihilation operator\hspace{0.9em}湮灭算子\par
ladder operator\hspace{0.9em}阶梯算子\par
ground state\hspace{0.9em}基态\par
excited state\hspace{0.9em}激发态\par
coherent state\hspace{0.9em}相干态\par
Feynman path integral\hspace{0.9em}费曼路径积分\par
time evolution operator\hspace{0.9em}时间演化算子\par
Schrödinger equation\hspace{0.9em}薛定谔方程\par
\end{multicols}

\filbreak
\vspace{0.9em}
\noindent\textbf{四、伪微分算子、象征与量子化}\par\nobreak\vspace{0.35em}
\begin{multicols}{2}
\raggedright
\setlength{\parindent}{0pt}
\setlength{\parskip}{0.16em}
\linespread{1.02}\selectfont
pseudodifferential operator\hspace{0.9em}伪微分算子\par
quantization\hspace{0.9em}量子化\par
quantization map\hspace{0.9em}量子化映射\par
symbol\hspace{0.9em}象征\par
principal symbol\hspace{0.9em}主象征\par
symbol class\hspace{0.9em}象征类\par
order function\hspace{0.9em}阶函数\par
order\hspace{0.9em}阶\par
elliptic\hspace{0.9em}椭圆的\par
hypoelliptic\hspace{0.9em}亚椭圆的\par
parametrix\hspace{0.9em}拟基本解\par
cutoff function\hspace{0.9em}截断函数\par
compactly supported\hspace{0.9em}紧支的\par
microlocal\hspace{0.9em}微局部的\par
wave front set\hspace{0.9em}波前集\par
support\hspace{0.9em}支集\par
singular support\hspace{0.9em}奇支集\par
Weyl quantization\hspace{0.9em}外尔量子化\par
Weyl ordering\hspace{0.9em}外尔排序\par
Weyl algebra\hspace{0.9em}外尔代数\par
normal ordering\hspace{0.9em}正规排序\par
symmetrization\hspace{0.9em}对称化\par
Wigner transform\hspace{0.9em}维格纳变换\par
Wigner function\hspace{0.9em}维格纳函数\par
Wigner distribution\hspace{0.9em}维格纳分布\par
Moyal product\hspace{0.9em}Moyal 乘积\par
star product\hspace{0.9em}星乘积\par
composition formula\hspace{0.9em}复合公式\par
Calderón–Vaillancourt\hspace{0.9em}卡尔德隆–韦扬库尔\par
\end{multicols}

\filbreak
\vspace{0.9em}
\noindent\textbf{五、几何量子化}\par\nobreak\vspace{0.35em}
\begin{multicols}{2}
\raggedright
\setlength{\parindent}{0pt}
\setlength{\parskip}{0.16em}
\linespread{1.02}\selectfont
symplectic\hspace{0.9em}辛的\par
symplectic manifold\hspace{0.9em}辛流形\par
symplectic structure\hspace{0.9em}辛结构\par
symplectic form\hspace{0.9em}辛形式\par
symplectic geometry\hspace{0.9em}辛几何\par
symplectomorphism\hspace{0.9em}辛同胚\par
Hamiltonian vector field\hspace{0.9em}哈密顿向量场\par
Poisson manifold\hspace{0.9em}泊松流形\par
Poisson structure\hspace{0.9em}泊松结构\par
Poisson algebra\hspace{0.9em}泊松代数\par
Lagrangian submanifold\hspace{0.9em}拉格朗日子流形\par
isotropic\hspace{0.9em}迷向的\par
polarization\hspace{0.9em}极化\par
real polarization\hspace{0.9em}实极化\par
Kähler polarization\hspace{0.9em}凯勒极化\par
prequantization\hspace{0.9em}预量子化\par
prequantum line bundle\hspace{0.9em}预量子化线丛\par
line bundle\hspace{0.9em}线丛\par
complex line bundle\hspace{0.9em}复线丛\par
connection\hspace{0.9em}联络\par
curvature\hspace{0.9em}曲率\par
Chern class\hspace{0.9em}陈类\par
integrality condition\hspace{0.9em}整性条件\par
geometric quantization\hspace{0.9em}几何量子化\par
geometric semi-classical quantization\hspace{0.9em}几何半经典量子化\par
Heisenberg group\hspace{0.9em}海森堡群\par
quantization condition\hspace{0.9em}量子化条件\par
Kähler potential\hspace{0.9em}凯勒势\par
complex structure\hspace{0.9em}复结构\par
holomorphic\hspace{0.9em}全纯的\par
leaf\hspace{0.9em}叶\par
group action\hspace{0.9em}群作用\par
orbit\hspace{0.9em}轨道\par
quotient\hspace{0.9em}商\par
reduction\hspace{0.9em}约化\par
differential form\hspace{0.9em}微分形式\par
Liouville measure\hspace{0.9em}刘维尔测度\par
Bohr–Sommerfeld condition\hspace{0.9em}玻尔–索末菲条件\par
\end{multicols}

\filbreak
\vspace{0.9em}
\noindent\textbf{六、形变量子化}\par\nobreak\vspace{0.35em}
\begin{multicols}{2}
\raggedright
\setlength{\parindent}{0pt}
\setlength{\parskip}{0.16em}
\linespread{1.02}\selectfont
deformation quantization\hspace{0.9em}形变量子化\par
formal deformation\hspace{0.9em}形式形变\par
formal power series\hspace{0.9em}形式幂级数\par
associative algebra\hspace{0.9em}结合代数\par
associativity\hspace{0.9em}结合性\par
Hochschild cohomology\hspace{0.9em}霍赫希尔德上同调\par
Maurer–Cartan equation\hspace{0.9em}毛雷尔–嘉当方程\par
formality\hspace{0.9em}形式性\par
Kontsevich\hspace{0.9em}孔采维奇\par
Fedosov\hspace{0.9em}费多谢夫\par
trace map\hspace{0.9em}迹映射\par
closed star product\hspace{0.9em}闭星乘积\par
trace functional\hspace{0.9em}迹泛函\par
character\hspace{0.9em}特征标\par
index theorem\hspace{0.9em}指标定理\par
Atiyah–Singer\hspace{0.9em}阿蒂亚–辛格\par
characteristic class\hspace{0.9em}示性类\par
Todd class\hspace{0.9em}托德类\par
noncommutative geometry\hspace{0.9em}非交换几何\par
bidifferential operator\hspace{0.9em}双微分算子\par
deformation parameter\hspace{0.9em}形变参数\par
Fréchet space\hspace{0.9em}弗雷歇空间\par
cohomology\hspace{0.9em}上同调\par
cocycle\hspace{0.9em}上闭链\par
obstruction\hspace{0.9em}障碍\par
\end{multicols}

\filbreak
\vspace{0.9em}
\noindent\textbf{七、数学通用词}\par\nobreak\vspace{0.35em}
\begin{multicols}{2}
\raggedright
\setlength{\parindent}{0pt}
\setlength{\parskip}{0.16em}
\linespread{1.02}\selectfont
manifold\hspace{0.9em}流形\par
chart\hspace{0.9em}图卡\par
coordinate\hspace{0.9em}坐标\par
coordinate system\hspace{0.9em}坐标系\par
tangent bundle\hspace{0.9em}切丛\par
cotangent bundle\hspace{0.9em}余切丛\par
section\hspace{0.9em}截面\par
vector field\hspace{0.9em}向量场\par
differential\hspace{0.9em}微分\par
closed\hspace{0.9em}闭的\par
exact\hspace{0.9em}恰当的\par
de Rham cohomology\hspace{0.9em}德拉姆上同调\par
Lie derivative\hspace{0.9em}李导数\par
interior product\hspace{0.9em}内乘\par
wedge product\hspace{0.9em}楔积\par
gradient\hspace{0.9em}梯度\par
Hessian\hspace{0.9em}黑塞\par
critical point\hspace{0.9em}临界点\par
nondegenerate\hspace{0.9em}非退化的\par
Morse\hspace{0.9em}莫尔斯\par
degree\hspace{0.9em}次数\par
homotopy\hspace{0.9em}同伦\par
homology\hspace{0.9em}同调\par
homogeneous space\hspace{0.9em}齐性空间\par
stabilizer\hspace{0.9em}稳定化子\par
bilinear form\hspace{0.9em}双线性形式\par
symmetric\hspace{0.9em}对称的\par
determinant\hspace{0.9em}行列式\par
inverse\hspace{0.9em}逆\par
map\hspace{0.9em}映射\par
isomorphism\hspace{0.9em}同构\par
diffeomorphism\hspace{0.9em}微分同胚\par
embedding\hspace{0.9em}嵌入\par
submanifold\hspace{0.9em}子流形\par
neighborhood\hspace{0.9em}邻域\par
open set\hspace{0.9em}开集\par
compact\hspace{0.9em}紧的\par
continuous\hspace{0.9em}连续的\par
smooth\hspace{0.9em}光滑的\par
analytic\hspace{0.9em}解析的\par
arbitrary\hspace{0.9em}任意的\par
unique\hspace{0.9em}唯一的\par
well-defined\hspace{0.9em}良定义的\par
trivial\hspace{0.9em}平凡的\par
nontrivial\hspace{0.9em}非平凡的\par
identically\hspace{0.9em}恒等地\par
up to\hspace{0.9em}相差\par
modulo\hspace{0.9em}模\par
respectively\hspace{0.9em}分别\par
equivalently\hspace{0.9em}等价地\par
in particular\hspace{0.9em}特别地\par
moreover\hspace{0.9em}此外\par
furthermore\hspace{0.9em}而且\par
note that\hspace{0.9em}注意\par
it follows that\hspace{0.9em}由此可知\par
as a result\hspace{0.9em}于是\par
hence\hspace{0.9em}从而\par
thus\hspace{0.9em}因此\par
indeed\hspace{0.9em}事实上\par
suppose\hspace{0.9em}设\par
assume\hspace{0.9em}假设\par
claim\hspace{0.9em}断言\par
satisfy\hspace{0.9em}满足\par
denote\hspace{0.9em}记\par
definite\hspace{0.9em}定的\par
bounded\hspace{0.9em}有界的\par
finite\hspace{0.9em}有限的\par
infinity\hspace{0.9em}无穷\par
sufficiently small\hspace{0.9em}充分小\par
sufficiently large\hspace{0.9em}充分大\par
leading order\hspace{0.9em}首阶\par
next order\hspace{0.9em}次阶\par
correction\hspace{0.9em}修正\par
approximation\hspace{0.9em}近似\par
exact solution\hspace{0.9em}精确解\par
closed form\hspace{0.9em}闭形式\par
explicit\hspace{0.9em}显式的\par
formula\hspace{0.9em}公式\par
expression\hspace{0.9em}表达式\par
inequality\hspace{0.9em}不等式\par
equality\hspace{0.9em}等式\par
equation\hspace{0.9em}方程（式）\par
system of equations\hspace{0.9em}方程组\par
condition\hspace{0.9em}条件\par
property\hspace{0.9em}性质\par
consequence\hspace{0.9em}推论\par
counterexample\hspace{0.9em}反例\par
hypothesis\hspace{0.9em}假设\par
proof\hspace{0.9em}证明\par
conclusion\hspace{0.9em}结论\par
remark\hspace{0.9em}注\par
definition\hspace{0.9em}定义\par
theorem\hspace{0.9em}定理\par
proposition\hspace{0.9em}命题\par
lemma\hspace{0.9em}引理\par
corollary\hspace{0.9em}推论\par
example\hspace{0.9em}例\par
notation\hspace{0.9em}记号\par
assumption\hspace{0.9em}假设\par
question\hspace{0.9em}问\par
problem\hspace{0.9em}问题\par
exercise\hspace{0.9em}练习\par
section\hspace{0.9em}节\par
chapter\hspace{0.9em}章\par
part\hspace{0.9em}部\par
appendix\hspace{0.9em}附录\par
bibliography\hspace{0.9em}参考文献\par
reference\hspace{0.9em}文献\par
index\hspace{0.9em}索引\par
preface\hspace{0.9em}前言\par
version\hspace{0.9em}版本\par
table of contents\hspace{0.9em}目录\par
figure\hspace{0.9em}图\par
stated\hspace{0.9em}所述的\par
\end{multicols}

\filbreak
\vspace{0.9em}
\noindent\textbf{八、人名（汉语定名；其余保留拉丁原文）}\par\nobreak\vspace{0.35em}
\begin{multicols}{2}
\raggedright
\setlength{\parindent}{0pt}
\setlength{\parskip}{0.16em}
\linespread{1.02}\selectfont
Hamilton\hspace{0.9em}哈密顿\par
Lagrange\hspace{0.9em}拉格朗日\par
Einstein\hspace{0.9em}爱因斯坦\par
Bohr\hspace{0.9em}玻尔\par
Dirac\hspace{0.9em}狄拉克\par
Schrödinger\hspace{0.9em}薛定谔\par
Heisenberg\hspace{0.9em}海森堡\par
Hilbert\hspace{0.9em}希尔伯特\par
Fourier\hspace{0.9em}傅里叶\par
Laplace\hspace{0.9em}拉普拉斯\par
Weyl\hspace{0.9em}外尔\par
Wigner\hspace{0.9em}维格纳\par
Feynman\hspace{0.9em}费曼\par
Planck\hspace{0.9em}普朗克\par
Kähler\hspace{0.9em}凯勒\par
Toeplitz\hspace{0.9em}托普利茨\par
Fredholm\hspace{0.9em}弗雷德霍姆\par
Stone\hspace{0.9em}斯通\par
von Neumann\hspace{0.9em}冯·诺依曼\par
Berezin\hspace{0.9em}别列津\par
Bargmann\hspace{0.9em}巴格曼\par
Fock\hspace{0.9em}福克\par
Schwartz\hspace{0.9em}施瓦茨\par
Stokes\hspace{0.9em}斯托克斯\par
Airy\hspace{0.9em}艾里\par
Maslov\hspace{0.9em}马斯洛夫\par
Hochschild\hspace{0.9em}霍赫希尔德\par
Maurer–Cartan\hspace{0.9em}毛雷尔–嘉当\par
Casimir\hspace{0.9em}卡西米尔\par
Liouville\hspace{0.9em}刘维尔\par
de Rham\hspace{0.9em}德拉姆\par
Todd\hspace{0.9em}托德\par
Fréchet\hspace{0.9em}弗雷歇\par
Moyal\hspace{0.9em}Moyal\par
Bargmann–Fock\hspace{0.9em}巴格曼–福克\par
Poisson\hspace{0.9em}泊松\par
Jacobi\hspace{0.9em}雅可比\par
Hamilton–Jacobi\hspace{0.9em}哈密顿–雅可比\par
Legendre\hspace{0.9em}勒让德\par
Hermite\hspace{0.9em}厄米\par
Gauss\hspace{0.9em}高斯\par
Wick\hspace{0.9em}威克\par
Wiener\hspace{0.9em}维纳\par
van Vleck\hspace{0.9em}范弗莱克\par
Kostant\hspace{0.9em}科斯坦特\par
Kohn–Nirenberg\hspace{0.9em}科恩–尼伦伯格\par
\end{multicols}

\filbreak
\vspace{0.9em}
\noindent\textbf{九 各章译名增补（终审合并）}\par\nobreak\vspace{0.35em}
\begin{multicols}{2}
\raggedright
\setlength{\parindent}{0pt}
\setlength{\parskip}{0.16em}
\linespread{1.02}\selectfont
separable (Hilbert space)\hspace{0.9em}可分的\par
Banach space\hspace{0.9em}巴拿赫空间\par
dominated convergence\hspace{0.9em}控制收敛定理\par
closable\hspace{0.9em}可闭的\par
closure (of an operator)\hspace{0.9em}闭包\par
extension (of an operator)\hspace{0.9em}扩张\par
deficiency index\hspace{0.9em}亏指数\par
locally constant\hspace{0.9em}局部常值\par
bounded below\hspace{0.9em}有下界\par
Cayley transform\hspace{0.9em}凯莱变换\par
isometry / isometric\hspace{0.9em}等距 / 等距映射\par
resolvent identity\hspace{0.9em}预解恒等式\par
analytic (operator-valued)\hspace{0.9em}解析的\par
Bergman kernel\hspace{0.9em}伯格曼核\par
heat operator\hspace{0.9em}热算子\par
shift operator\hspace{0.9em}移位算子\par
overcomplete family\hspace{0.9em}超完备族\par
multi-index\hspace{0.9em}多重指标\par
reproducing property\hspace{0.9em}再生性质\par
strict deformation quantization\hspace{0.9em}严格形变量化\par
de Broglie relations\hspace{0.9em}德布罗意关系\par
matter wave\hspace{0.9em}物质波\par
dispersion relation\hspace{0.9em}色散关系\par
group velocity\hspace{0.9em}群速度\par
phase velocity\hspace{0.9em}相速度\par
Gaussian packet\hspace{0.9em}高斯波包\par
stationary states\hspace{0.9em}定态\par
time-independent Schrödinger equation\hspace{0.9em}定态薛定谔方程\par
raising operator\hspace{0.9em}升算子\par
lowering operator\hspace{0.9em}降算子\par
Hermite polynomials\hspace{0.9em}厄米多项式\par
number operator\hspace{0.9em}数算子\par
probability current\hspace{0.9em}概率流\par
deep well limit\hspace{0.9em}深阱极限\par
resonant transmission\hspace{0.9em}共振透射\par
even function\hspace{0.9em}偶函数\par
odd function\hspace{0.9em}奇函数\par
Pauli matrices\hspace{0.9em}泡利矩阵\par
singlet state\hspace{0.9em}单态\par
Bell's inequality\hspace{0.9em}贝尔不等式\par
Bell\hspace{0.9em}贝尔\par
Casimir element\hspace{0.9em}卡西米尔元\par
hidden variable\hspace{0.9em}隐变量\par
local hidden variable theory\hspace{0.9em}定域隐变量理论\par
spooky action at a distance\hspace{0.9em}幽灵般的超距作用\par
Levi-Civita symbol\hspace{0.9em}Levi-Civita 符号\par
Alice / Bob\hspace{0.9em}Alice / Bob\par
essential range\hspace{0.9em}本质值域\par
delta function\hspace{0.9em}delta 函数\par
distribution space\hspace{0.9em}分布空间\par
cyclic subspace\hspace{0.9em}循环子空间\par
resolvent-invariant\hspace{0.9em}预解不变的\par
spectral basis\hspace{0.9em}谱基\par
Herglotz–Nevanlinna function\hspace{0.9em}赫格洛茨–涅万林纳函数\par
Herglotz representation\hspace{0.9em}赫格洛茨表示\par
Riesz representation theorem\hspace{0.9em}里斯表示定理\par
probability measure\hspace{0.9em}概率测度\par
characteristic function (of a set)\hspace{0.9em}指示函数\par
polarization identity\hspace{0.9em}极化恒等式\par
isolated point\hspace{0.9em}孤立点\par
accumulation point\hspace{0.9em}聚点\par
constant projection\hspace{0.9em}常投影\par
Weyl criterion\hspace{0.9em}外尔判别法\par
Weyl sequence\hspace{0.9em}外尔序列\par
uniform boundedness theorem\hspace{0.9em}一致有界性定理\par
two-sided ideal\hspace{0.9em}双边理想\par
finite rank\hspace{0.9em}有限秩\par
relatively compact\hspace{0.9em}相对紧\par
diagonalization argument\hspace{0.9em}对角线论证\par
Schwartz kernel\hspace{0.9em}施瓦茨核\par
Cauchy–Schwarz inequality\hspace{0.9em}柯西–施瓦茨不等式\par
compact perturbation\hspace{0.9em}紧扰动\par
infinitesimal generator\hspace{0.9em}无穷小生成元\par
Borel function calculus\hspace{0.9em}博雷尔函数演算\par
unitary equivalent\hspace{0.9em}酉等价\par
sesquilinear form\hspace{0.9em}半双线性形式\par
modified Bessel function\hspace{0.9em}修改的贝塞尔函数\par
Hermite polynomial\hspace{0.9em}厄米多项式\par
Kato's inequality\hspace{0.9em}加藤不等式\par
diagonal argument\hspace{0.9em}对角线论证\par
semi-bounded\hspace{0.9em}半有界的\par
relatively bounded\hspace{0.9em}相对有界\par
Kato–Rellich theorem\hspace{0.9em}加藤–雷利希定理\par
Rellich\hspace{0.9em}雷利希\par
Trotter product formula\hspace{0.9em}特罗特乘积公式\par
Lie product formula\hspace{0.9em}李乘积公式\par
graph norm\hspace{0.9em}图像范数\par
Five Lemma\hspace{0.9em}五引理\par
open mapping theorem\hspace{0.9em}开映射定理\par
Atkinson's theorem\hspace{0.9em}阿特金森定理\par
Atkinson\hspace{0.9em}阿特金森\par
Banach algebra\hspace{0.9em}巴拿赫代数\par
Calkin algebra\hspace{0.9em}卡尔金代数\par
Hardy space\hspace{0.9em}哈代空间\par
Toeplitz operator\hspace{0.9em}托普利茨算子\par
Toeplitz index theorem\hspace{0.9em}托普利茨指标定理\par
symplectic basis\hspace{0.9em}辛基\par
symplectomorphic (vector spaces)\hspace{0.9em}辛同构\par
indicator function (characteristic function of a set)\hspace{0.9em}示性函数\par
Bargmann–Fock representation\hspace{0.9em}巴格曼–福克表示\par
Segal–Bargmann transform\hspace{0.9em}Segal–Bargmann 变换\par
prequantization condition\hspace{0.9em}预量子化条件\par
quantizable\hspace{0.9em}可量子化的\par
Liouville 1-form\hspace{0.9em}刘维尔 1-形式\par
hermitian holomorphic line bundle\hspace{0.9em}厄米全纯线丛\par
positive line bundle\hspace{0.9em}正线丛\par
Euler characteristic\hspace{0.9em}欧拉示性数\par
Hirzebruch–Riemann–Roch theorem\hspace{0.9em}希尔策布鲁赫–黎曼–罗赫定理\par
half-density\hspace{0.9em}半密度\par
canonical line bundle\hspace{0.9em}典范线束\par
square root (of a line bundle)\hspace{0.9em}平方根\par
Dirichlet\hspace{0.9em}狄利克雷\par
Heaviside\hspace{0.9em}赫维赛德\par
Grassmann\hspace{0.9em}格拉斯曼\par
Gram–Schmidt\hspace{0.9em}格拉姆–施密特\par
Morette\hspace{0.9em}莫雷特\par
Rieffel\hspace{0.9em}里费尔\par
Baker–Campbell–Hausdorff\hspace{0.9em}贝克–坎贝尔–豪斯多夫\par
Keller\hspace{0.9em}凯勒\par
Kato–Rellich\hspace{0.9em}加藤–雷利希\par
DGA\hspace{0.9em}微分分次代数（DGA）\par
EBK\hspace{0.9em}爱因斯坦–布里渊–凯勒（EBK）\par
Groenewold–Van Hove theorem\hspace{0.9em}格罗嫩沃尔德–范霍夫定理\par
DeWilde–Lecomte\hspace{0.9em}德维尔德–勒孔特\par
Hochschild–Kostant–Rosenberg\hspace{0.9em}霍赫希尔德–科斯坦特–罗森伯格\par
Schwartz function\hspace{0.9em}施瓦茨函数\par
Poincaré–Birkhoff–Witt\hspace{0.9em}庞加莱–伯克霍夫–维特\par
Hugh Everett\hspace{0.9em}埃弗里特\par
David Mermin\hspace{0.9em}默明\par
potential function\hspace{0.9em}势函数\par
EPR\hspace{0.9em}爱因斯坦–波多尔斯基–罗森佯谬（EPR）\par
\end{multicols}

\endgroup
